\documentclass[book.tex]{subfiles}

\begin{document}
\chapter{Quantum Programming}

\label{chap:quantum_programmiing}
\noindent\rule{11cm}{0.4pt} \\
\textbf{Prerequisites:} \autoref{chap:physika} \\
\textbf{Difficulty Level:} ** \\
\noindent\rule{11cm}{0.4pt}
\vspace{5mm}


Quantum computing has spawned a large number of new programming languages to design sophisticated quantum algorithms. While quantum circuits are themselves a simple language for quantum algorithms, circuits can become verbose. Most quantum algorithms also have repeated patterns within circuits suggesting that there are higher order patterns which can be encapsulated and re-sued in circuit design.

In this section, we will discuss a simple extension to Physika that performs quantum programming. We will design this quantum extension to enable variational quantum programs which will require us to merge classical and quantum computations in the same program.

\section{Syntax for Quantum Gates}

Basic gates are provided by simple functions
\begin{lstlisting}
X : $H_2 \to H_2$
Y : $H_2 \to H_2$
$\sigma_x$ : $H_2 \to H_2$
\end{lstlisting}

These operators can be chained together into basic algorithms and executed the same as any other algorithm. The key differentiator is that the type must be from a specified Hilbert space to another specified Hilbert space.
%\textcolor{red}{TODO: Add in projectors / SWAP test}

%\textcolor{red}{TODO: Add in syntax for Pauli strings}

%\section{Parameterized Quantum Circuits}

%\textcolor{red}{TODO: Add in syntax for parameterized quantum gates}

\end{document}
